\section{A weighted Gaussian mixture}
\label{app:weighted-mixture}

We give an independent Gaussian construction for the mean-cost
argument of Section~\ref{sec:mean-cost}. The event again requires
exactly the selected leaf to fail its interval test, but now we first
divide the central coordinate into bins and choose each leaf interval
from its bin's conditional mean. This lets us compute the variance on
the event as a finite mixture. The relevant weights change when we
condition on the event; controlling those weights is the main estimate.

Our goal is a positive-probability event on which the sum has
conditional variance less than one, with an indicator that saves one
coordinate in each summand. Lemma~\ref{lem:refinement} will then
produce the required finite partition. The parameter choice below is
independent of the parameter choice in Lemma~\ref{lem:gaussian}.

\subsection{The event and its conditional bin weights}
\label{app:weighted-identity}

For $R>0$, put
\begin{equation}\label{app:weighted-schedule}
 m=1+\lceil e^{R^2/4}\rceil,\qquad
 \Delta=\frac{2R}{m-1},\qquad h=e^{-3R^2/8},\qquad q=m+1.
\end{equation}
Let $G=(G_0,\ldots,G_m)$ have independent standard normal coordinates
and set $L=\sum_{j=0}^mG_j$. Partition $[-R,R)$ into $m-1$ half-open
intervals $B_i$ of length $\Delta$, and let
$B_m=\R\setminus[-R,R)$. Define
\[
 \pi_i=\Pp(G_0\in B_i),\qquad
 v_i=\E[G_0\mid G_0\in B_i],\qquad
 u_i=\Var(G_0\mid G_0\in B_i).
\]
Every $\pi_i$ is positive. For each leaf $j\in\{1,\ldots,m\}$,
take the half-open interval
\[
 I_j=[v_j+1-h/2,v_j+1+h/2),\qquad \chi_j=\one_{I_j}.
\]
For any real input $g=(g_0,\ldots,g_m)$, define $g\in A$ when,
for the unique $i$ such that $g_0\in B_i$, all leaves other than $i$
belong to their intervals and leaf $i$ does not. The half-open
conventions assign every endpoint a definite outcome.
All unqualified bin and leaf sums below range from $1$ to $m$.

\begin{lemma}[The weighted-mixture event]\label{lem:weighted-event}
For all sufficiently large finite $R$, the event $A$ just defined
satisfies
\[
 p_A:=\Pp(G\in A)>0,\qquad r_A:=\Var(L\mid G\in A)<1.
\]
Its boundary is contained in finitely many coordinate hyperplanes,
and on every real input its indicator has the expansion
\begin{equation}\label{app:weighted-cell}
 \one_A(g)=
 \sum_{i=1}^m\one_{B_i}(g_0)
       \prod_{\substack{1\le j\le m\\j\ne i}}\chi_j(g_j)
       -\prod_{j=1}^m\chi_j(g_j).
\end{equation}
Each product involves at most $m=q-1$ coordinates.
\end{lemma}

\begin{proof}
Exactly one bin contains $g_0$. The first term in
\eqref{app:weighted-cell} tests all unselected leaves; subtracting
the all-leaves-pass indicator requires the selected leaf to fail.
This proves the identity pointwise, including all endpoints.
The boundary assertion follows from the finite collections of bin
and leaf-interval endpoints. It remains to prove the variance saving.
We let $R\to\infty$ for this calculation and fix a finite value only
after the estimates.

\paragraph{Conditional moments.}
Write $\phi$ for the standard normal density. Symmetry of the two
tails gives $v_m=0$; for $i<m$, the conditional mean lies in the
closure of $B_i$, so $|v_i|\le R$ and $u_i\le\Delta^2$.
Integration by parts gives
\[
 \int_R^\infty t^2\phi(t)\,dt
 =R\phi(R)+\int_R^\infty\phi(t)\,dt,
 \qquad
 \int_R^\infty\phi(t)\,dt\le\frac{\phi(R)}R.
\]
Thus
\begin{equation}\label{app:weighted-bin-moments}
 \pi_m u_m=O(Re^{-R^2/2}),\qquad
 \sum_i\pi_i v_i=0,\qquad
 \sum_i\pi_i v_i^2=1-\sum_i\pi_i u_i=1-o(1),
\end{equation}
where the last equality follows from
$\sum_i\pi_i u_i\le\Delta^2+O(Re^{-R^2/2})=o(1)$.

For the leaf intervals, write
\[
 \alpha_j=\Pp(G_j\in I_j),\qquad
 t_j=\E[G_j\mid G_j\in I_j],\qquad
 r_j=\Var(G_j\mid G_j\in I_j).
\]
The bounded normal density and the interval lengths imply, uniformly
in $j$ for sufficiently large $R$,
\begin{equation}\label{app:weighted-rare-moments}
 0<\alpha_j\le Ch<\tfrac12,\qquad
 |t_j-(v_j+1)|\le h/2,\qquad 0\le r_j\le h^2.
\end{equation}
Constants in this section are absolute unless stated otherwise.
Let $d_j,b_j$ be the conditional mean and variance of $G_j$ on
$I_j^c$. Using the unconditional mean zero and variance one gives
\begin{equation}\label{app:weighted-complement}
 d_j=-\frac{\alpha_jt_j}{1-\alpha_j},\qquad
 b_j=\frac{1-\alpha_jr_j}{1-\alpha_j}
       -\frac{\alpha_jt_j^2}{(1-\alpha_j)^2}.
\end{equation}
Indeed the complementary second moment is
$(1-\alpha_j(r_j+t_j^2))/(1-\alpha_j)$, from which we subtract
$d_j^2$.

\paragraph{The mixture identity.}
The event $A$ changes the distribution of the selected bin. Put
\[
 w_i=\pi_i\frac{1-\alpha_i}{\alpha_i},\qquad
 W=\sum_iw_i,\qquad \omega_i=\frac{w_i}{W}.
\]
Independence gives
\begin{equation}\label{app:weighted-probability}
 p_A=W\prod_{j=1}^m\alpha_j>0,
 \qquad \Pp(G_0\in B_i\mid G\in A)=\omega_i.
\end{equation}
Here $W$ is finite and positive for every finite $R$. Within
$A\cap\{G_0\in B_i\}$ the coordinates remain independent under
their respective bin, interval, and complementary-interval restrictions.
The conditional mean and variance of $L$ are therefore
\[
 \sum_{j=1}^mt_j-1+e_i,
 \qquad u_i+\sum_jr_j+b_i-r_i,
 \qquad e_i=v_i+1-\frac{t_i}{1-\alpha_i}.
\]
The small mean error satisfies
\begin{equation}\label{app:weighted-mean-error}
 |e_i|\le h/2+\frac{\alpha_i|t_i|}{1-\alpha_i}
          =O(h(R+2)).
\end{equation}
Applying the variance decomposition with the conditional weights
$\omega_i$, not the original weights $\pi_i$, gives the exact identity
\begin{align}
 W\bigl(1-\Var(L\mid G\in A)\bigr)
 &=\sum_iw_i(1-b_i+r_i)-\sum_iw_iu_i\notag\\
 &\quad-W\sum_jr_j-W\Var_{i\sim\omega}(e_i).
 \label{app:weighted-variance-identity}
\end{align}
We next show that the first sum has lower limit at least one and
that each of the three subtracted terms tends to zero. This will
give the strict conditional variance bound despite the possibly
large total weight $W$.

The schedule in \eqref{app:weighted-schedule} makes the bin scale
$\Delta^2/h$, the aggregate leaf scale $mh$, and the tail scale
$e^{-R^2/2}/h$ tend to zero even after multiplication by $e^{3R}$
and any fixed power of $R$. The leaf intervals are narrow enough to
control their total variance, but wide enough that inverse interval
probabilities do not overwhelm the bin and tail estimates below.

\paragraph{The positive contribution.}
Substituting \eqref{app:weighted-complement} yields
\begin{equation}\label{app:weighted-positive-term}
 w_i(1-b_i+r_i)
 =\pi_i\left(-1+\frac{r_i}{\alpha_i}
                    +\frac{t_i^2}{1-\alpha_i}\right)
 \ge\pi_i(t_i^2-1).
\end{equation}
Since $t_i=v_i+1+O(h)$ uniformly, we have
\[
 \sum_i\pi_i(t_i^2-1)
 =\sum_i\pi_i v_i^2+2\sum_i\pi_i v_i+O(h(R+2))
 =1+o(1)
\]
by \eqref{app:weighted-bin-moments}. The shift of the leaf intervals
by one is important: it produces the term $t_i^2-1$, whose average
tends to one.

\paragraph{The three error terms.}
For $i<m$, the fact that $v_i$ belongs to the closure of $B_i$
gives the density-ratio bound
\begin{align}
 w_i\le\frac{\pi_i}{\alpha_i}
 &\le\frac{\Delta}{h}
       \frac{\phi(\max\{0,|v_i|-\Delta\})}
            {\phi(|v_i|+1+h)}\notag\\
 &\le\frac{\Delta}{h}e^{3R}
 \label{app:weighted-density-ratio}
\end{align}
for all sufficiently large $R$. For the final inequality, put
$x=|v_i|\le R$. The logarithm of the density ratio is at most
\[
 \frac{(x+1+h)^2-\max\{0,x-\Delta\}^2}{2}
 \le(1+h+\Delta)x+\frac{(1+h)^2}{2}\le3R.
\]
The middle inequality follows by expanding when $x\ge\Delta$;
when $x<\Delta$, use $x^2/2\le\Delta x$.
For the tail bin, $v_m=0$, so
$\alpha_m\ge h\phi(1+h)$ and $w_m\le C\pi_m/h$.
Summing and using $(m-1)\Delta=2R$ gives
\begin{equation}\label{app:weighted-total-weight}
 W\le\frac{2Re^{3R}+C}{h}.
\end{equation}

First, $u_i\le\Delta^2$ on interior bins and
$\pi_mu_m=O(Re^{-R^2/2})$ on the tail bin. Since
$\Delta^2\le4R^2e^{-R^2/2}$, we obtain
\[
 \sum_iw_iu_i
 \le W\Delta^2+C\pi_mu_m/h
 =O(R^3e^{3R-R^2/8})+O(Re^{-R^2/8})=o(1).
\]
Second, $\sum_jr_j\le mh^2$ and $mh=O(e^{-R^2/8})$, so
\[
 W\sum_jr_j\le Wmh^2
 =O((Re^{3R}+1)e^{-R^2/8})=o(1).
\]
Third, \eqref{app:weighted-mean-error} gives
\[
 W\Var_{i\sim\omega}(e_i)
 \le W\,O(h^2(R+2)^2)
 =O((Re^{3R}+1)(R+2)^2e^{-3R^2/8})=o(1).
\]
Each negative quadratic exponent dominates the linear exponent and
polynomial factors. Combining these estimates with
\eqref{app:weighted-variance-identity} and
\eqref{app:weighted-positive-term} shows that
\[
 W\bigl(1-\Var(L\mid G\in A)\bigr)\ge1+o(1)>0
\]
for all sufficiently large $R$. Since $W>0$, this proves the lemma.
\end{proof}

\subsection{Application to mean-cost amplification}
\label{app:weighted-transfer}

Fix one finite $R$ for which Lemma~\ref{lem:weighted-event} holds.
Thus the integer $q=m+1$, all intervals and bins, and the positive
numbers
\begin{equation}\label{app:weighted-saving}
 p_A=\Pp(G\in A)>0,\qquad 1-r_A>0
\end{equation}
are fixed before further approximation. The pointwise expansion
\eqref{app:weighted-cell} uses at most $q-1$ coordinates in each
term, and the finite hyperplane boundary is Gaussian-null. These
are precisely the hypotheses of Lemma~\ref{lem:refinement}.
That lemma refines $A^c$ into a finite collection of cells, with
total squared approximation error below the actual positive quantity
$p_A(1-r_A)$. It assigns cost $m$ to $A$ and cost $q$ to every
other cell, giving a fixed partition $\mathcal P$ for which
\begin{equation}\label{app:weighted-mean-cost-gain}
 V:=\Var(\E[L\mid\mathcal P(G)])>
 c:=\E\ell(\mathcal P(G))=q-p_A.
\end{equation}
The partition includes its Gaussian-null cells as pointwise sets;
later discrete input laws may give those cells positive probability.

With $R$ and this partition fixed, Lemma~\ref{lem:mc-transfer}
supplies the finite-input transfer used in Theorem~\ref{thm:mean-cost}.
The costs $m$ and $q$ are positive integers, and the pointwise coordinate
bounds, Gaussian-null boundaries, and strict inequality $V>c$ verify
all its hypotheses. The theorem therefore gives arbitrarily large
ratios of retained variance to cell degree and the signed Boolean conclusion
of Theorem~\ref{thm:main}. Neither the conditional variance saving nor
the required batch sizes are asserted to be uniform in $R$.
